\documentclass[book.tex]{subfiles}

\begin{document}

\chapter{Path Integral Theory of an Oscillator}

\label{chap:perturbation_theory}
\noindent\rule{11cm}{0.4pt} \\
\textbf{Prerequisites:} \autoref{chap:schrodinger} \\
\textbf{Difficulty Level:} ** \\
\noindent\rule{11cm}{0.4pt}
\vspace{5mm}

%\textcolor{red}{TODO: Cite Wipf}

The action of a harmonic oscillator with mass $m$ is %(\textcolor{red}{TODO: Do a chapter on the classical action for a harmonic oscillator.} 
\begin{align}
    S &= \frac{m}{2} \int_{t'}^{t} ds(\dot{w}^2(s)-\omega^2(s) w^2(s))
\end{align}
Here $\omega(s)$ is a time-dependent circular frequency. Consider the following path
\begin{align}
S^{(n)}(w) &= \frac{m}{2}\sum_{j=0}^{n-1}\left [ \frac{1}{\epsilon}(w_{j+1}-w_j)^2 - \epsilon \omega_j^2 w_j^2 \right ] \\
w_{j+1} &= \omega (t' + j\epsilon)
\end{align}

\begin{align}
\xi &= (w_1,\dotsc,,w_{n-1}) \\
\eta &= (q', 0,\dotsc,0, q)
\end{align}
Then we can write action as
\begin{align}
S^{(n)}(w) &= \frac{m}{2} \left ( \frac{1}{\epsilon} (\eta, \eta)  + \frac{1}{\epsilon} (\xi, C \xi) - \frac{2}{\epsilon} (\xi, \epsilon) - \epsilon \omega_0^2 (q')^2 \right )
\end{align}
Here $C$ is a $n-1$ dimensional matrix. The action then reduces to computing a Gaussian integral on a discrete time lattice. The derivation yields inthe end
\begin{align}
K(t, q, t', q') &= \lim_{\epsilon \to 0} \sqrt{\frac{m}{2\pi i \hbar}} \frac{1}{\epsilon \det C} e^{iS(n)(\xi)/\hbar}
\end{align}

\end{document}